\documentclass{article}
\usepackage{xeCJK}
\setCJKmainfont{Malgun Gothic}
\usepackage{geometry}
\usepackage{amsmath}
\usepackage{listings}
\usepackage{xcolor}

%\usepackage{kotex}

\usepackage{amsmath}
\usepackage{amssymb}
\usepackage{bm}
\usepackage{mathrsfs}
\usepackage{eucal}
\usepackage{mathtools}
\usepackage{unicode-math}

\lstdefinelanguage{CSharp}
{
morekeywords={public,static,void,int,float,string,return,class,if,else,new,var,delegate,Func,bool,using,namespace,partial,get,set,private,params,in},
sensitive=true,
morecomment=[l]{//},
morecomment=[s]{/*}{*/},
morestring=[b]",
}
\lstset{
  language=CSharp,
  basicstyle=\ttfamily\small,     % 기본 폰트
  keywordstyle=\color{blue}\bfseries,     % 키워드 색
  commentstyle=\color{gray},              % 주석 색
  stringstyle=\color{red},                % 문자열 색
  numbers=left,                           % 줄 번호
  numberstyle=\small\color{gray},
  stepnumber=1,
  numbersep=10pt,
  backgroundcolor=\color{white},          % 배경색
  frame=single,                           % 테두리
  tabsize=2,
  captionpos=b,
  breaklines=true,
  breakatwhitespace=false,
  showspaces=false,
  showstringspaces=false
}

\geometry{
 a4paper,
 left=10mm,
 right=10mm,
 top=10mm,
 bottom=20mm
 }

\title{재귀\,함수에\,관한\,고찰}
\author{사서림}
\date{2025.04.13.}

\setlength{\parindent}{0pt}

\begin{document}
\maketitle

\newpage
[재귀\,함수에\,관한\,고찰]
\\\\
재귀\,함수는\,어떤\,함수에\,대해서\,조건이\,부합할\,때까지만\,자기자신,\,즉\,지금\,실행하고\,있는\,함수를\,실행시키는\,함수이다.\\
우리는\,이\,함수에서\,두\,가지\,요소를\,쉽게\,발견\,할\,수\,있다.\,
바로\,실행할\,함수와\,계속\,반복\,실행할\,조건이다.
\begin{center}
$f=$실행할\,함수
\\$g=$조건
\end{center}
\,

이렇게\,표기하고\,재귀\,함수를\,자세히\,들여다보면\,우리는\,함수적\,표현으로\,재귀\,함수를\,$R$라고\,표현했을\,때\,다음과\,같이\,써볼\,수도\,있다는\,생각이\,들\,것이다.
\begin{center}
$R(f,g)(x)$
\end{center}
\,

그렇다면\,재귀\,함수\,$R(f,g)(x)$의\,공역은\,무엇일까?\,한번\,자세히\,생각해보자.\\
일단\,우리는\,재귀\,함수는\,조건에\,부합할\,때까지만\,실행\,중인\,$f$에서\,$f$안에\,$f$를\,넣는다는\,걸\,알\,수\,있다.\,
이때\,함수에\,넣는\,값도\,보기위하여\,함수의\,표현을\,써서\,$f(x)$라고\,하면\,우리는\,다음과\,같음을\,알\,수\,있다.
\begin{center}
조건에\,부합하면\,$f(x)$실행\,중에\,$f(x)$에\,$f$를\,넣는다.\,즉\\
$g=$\,true\,$\Rightarrow f(f)$
\end{center}
\,

즉\,$f$의\,공역은\,$f$의\,정의역과\,같아야\,함을\,알\,수\,있다.\,우리는\,이것을\,이렇게\,쓸\,수\,있다.
\begin{center}
$f:T \rightarrow T$
\end{center}
\,

이때\,우리는\,$g$또한\,함수로\,다룰\,수\,없을까\,생각할\,수\,있다.\,
그럼\,$g$의\,정의역은\,무엇일까?\,
우리는\,보통\,조건을\,$f$의\,입력값\,$x$에\,대해서\,생각한다.\,
하지만\,조건이\,$f(x)$에\,영향을\,받는\,특이한\,경우도\,있으므로\,$g$를\,다음과\,같이\,정의할\,수\,있다.
\begin{center}
$g(x,f)$\\
$g:T \times (T \rightarrow T) \rightarrow$\,bool
\end{center}
\,

그럼\,$R(f,g)(x)$는\,다음과\,같이\,표기\,가능하다.
\begin{center}
$R:T \rightarrow T$\\
$R(f,g)(x)=
\begin{cases}
x & \text{if}\,g(x,f(x))=\text{true} \\
R(f,g)(f(x)) & \text{otherwise}
\end{cases}
$
\end{center}
\,

하지만\,우리\,한번\,굳이\,조건이\,끝났을\,때\,무조건\,입력값\,$x$를\,반환해야\,하는지에\,대해서\,생각해보자.\,
사실\,잘\,생각해보면\,조건이\,끝났을\,때\,어떻게\,연산(동작)하고\,끝낼지도\,결정할\,수\,있다.\,이를\,$h$라고\,해보자.\,
$h$는\,$x,\,f,\,g$의\,영향을\,받을\,수\,있는데\,이는\,입력값$x$와\,두\,개의\,함수\,$f,\,g$가\,매개변수로서\,동작해야\,한다는\,말로도\,이해할\,수\,있다.\,
그러므로\,$h$를\,다음과\,같이\,정의하자.
\begin{center}
$h(x,f,g)$
\end{center}
\,

그런대\,여기에서\,잘\,생각해보면\,꼭\,$h$의\,공역이\,$f$의\,정의역이나\,공역과\,같을\,필요는\,없다.\,
왜냐하면\,$h(x,f,g)$는\,$f(x)$를\,매개변수로\,받을\,뿐\,순수하게\,독립된\,함수이기\,때문이다.\,
그러므로\,우리는\,$R$의\,공역을\,$f$의\,정의역이나\,공역인\,$T$와\,다르게\,다른\,집합\,$U$로\,잡을\,수도\,있다.\,
그러므로\,$h$를\,정의하고\,$R$을\,다시\,정의할\,수\,있다.
\begin{center}
$f(x)=\text{함수}$\\
$g(x,f(x))=\text{함수}f(x)\text{의\,재귀가\,종료될\,조건}$\\
$h(x,f(x),g(x,f(x)))=\text{재귀가\,끝난\,후\,연산할\,함수(즉\,해석기)}$
\end{center}
\,

여기에서\,각각\,함수의\,정의는\,다음과\,같다.
\begin{center}
$f:T \rightarrow T$\\
$g:T \times (T \rightarrow T) \rightarrow \text{bool}$\\
$h:T \times (T \rightarrow T) \times (T \times (T \rightarrow T) \rightarrow \text{bool}) \rightarrow U$
\end{center}
\,

그리고\,최종적인\,재귀\,함수\,$R$은\,다음과\,같이\,정의한다.
\begin{center}
$R:f \times g \times h \times T \rightarrow U$\\
$R:(T \rightarrow T) \times (T \times (T \rightarrow T) \rightarrow \text{bool}) \times (T \times (T \rightarrow T) \times (T \times (T \rightarrow T) \rightarrow \text{bool}) \rightarrow U) \times T \rightarrow U$\\
$\text{재귀의\,핵심\,입력과\,출력간의\,관계만\,따지면\,}R:T \rightarrow U$\\
$R(f,g,h)(x)=
\begin{cases}
h(x,f,g) &
\text{if}\,g(x,f(x))=\text{true} \\
R(f,g,h)(f(x)) &
\text{otherwise}
\end{cases}
$
\end{center}
\,
\newpage

재귀\,함수$R(f,g,h)(x)$는\,기존에\,통합되어\,있던\,실제\,동작하는\,함수와\,반복\,조건\,그리고\,결과에\,대한\,연산부를\,독립적인\,함수로\,분리함으로써\,이전보다\,편하게\,유지관리할\,수\,있으며\,재귀\,함수\,자체를\,면밀히\,분석하여\,기존보다\,더\,잘\,관리할\,수\,있는\,가능성을\,제시한다.\,
그리고\,재귀\,함수$R$을\,변수로도\,볼\,수\,있을것이라는\,생각\,또한\,적어본다.\\
\\
-[재귀\,함수에\,대한\,고찰]끝-
\newline
\newline
\newline
\newline

[재귀\,함수\,R을\,C\#으로\,구현\,시도해보기]
\\\\
이번\,챕터에서는\,C\#를\,이용하여\,$R(f,g,h)(x)$을\,구현하는\,것을\,시도\,해보고자\,한다.\,
그럼\,바로\,클래스를\,만들어서\,필요한\,변수부터\,정의하도록\,하자.
\\
\newline
\begin{lstlisting}[caption={재귀\,함수\,R의\,기본\,매개변수\,$f,g,h$에\,대한\,정의부}]
namespace RecursionFuntion
{
    public partial class RecursionFuntion<T,U>
    {
        public Func<T, T> f { get; private set; }
        public Func<T,Func<T,T>, bool> g { get; private set; }
        public Func<T,Func<T,T>, Func<T, Func<T, T>, bool>,U> h { get; private set; }
    }
}
\end{lstlisting}

위의\,내용을\,보면\,RecursionFuntion클래스를\,템플릿\,T,\,U를\,사용하는\,클래스로\,만들어서\,$R:T \rightarrow U$개념을\,구현하려고\,하였음을\,알\,수\,있다.\,
그리고\,각\,변수\,$f,g,h$를\,정의한\,것을\,보면\,충실하게\,아래의\,구조를\,잘\,따르게\,만들어져\,있음을\,알\,수\,있다.\\
\newline
$f:T \rightarrow T$ \,\,\,\, Func<매개변수,출력>임으로\\Func<T,T>로\,정의됨 \\
\newline
$g:T \times (T \rightarrow T) \rightarrow \text{bool}$ \,\,\,\,  Func<매개변수1,매개변수2,출력>임으로\\Func<T,Func<T,T>, bool>로\,정의됨 \\
\newline
$h:T \times (T \rightarrow T) \times (T \times (T \rightarrow T) \rightarrow \text{bool}) \rightarrow U$ \,\,\,\, Func<매개변수1,매개변수2,매개변수3,출력>임으로\\Func<T,Func<T,T>, Func<T, Func<T, T>, bool>,U>로\,정의됨 \\
\newline

여기에서\,보면\,$f$에서\,$g,\,h$로\,갈수록\,이전\,함수에\,의존적이라는것을\,알\,수\,있다.\,
이는\,$f(x),\,g(x,f),\,h(x,f,g)$로서\,조건$g$가\,$f$의\,영향을\,받고,\,해석기h는\,$f$와\,$g$의\,영향을\,받을\,수도\,있기\,때문이다.\\
이제\,생성자와\,$f,g,h$를\,햘당할\,수\,있는\,메소드를\,만들어보자.
\\
\newline
\begin{lstlisting}[caption={재귀\,함수\,R의\,생성자\,및\,햘당\,메소드}]
namespace RecursionFuntion
{
    public partial class RecursionFuntion<T,U>
    {
        public RecursionFuntion() { }
        public RecursionFuntion(
            Func<T, T> f,
            Func<T, Func<T, T>, bool> g,
            Func<T, Func<T, T>, Func<T, Func<T, T>, bool>, U> h)
        {
            this.f = f;
            this.g = g;
            this.h = h;
        }
        public void SetF(Func<T, T> f)
        {
            this.f = f;
        }
        public void SetG(Func<T, Func<T, T>, bool> g)
        {
            this.g = g;
        }
        public void SetH(Func<T, Func<T, T>, Func<T, Func<T, T>, bool>, U> h)
        {
            this.h = h;
        }
    }
}
\end{lstlisting}

위의\,내용을\,보면\,말\,그대로\,$f,g,h$를\,햘당하는\,간단한\,생성자와\,메소드들이\,있다.\,
앞으로\,이것을\,어떻게\,사용할\,수\,있는지는\,실제로\,팩토리얼\,재귀\,함수를\,만들어서\,확인해보겠다.\,
이제\,R\,재귀\,함수의\,구동부\,메소드를\,보자.\,
쉬운\,이해를\,위하여\,선형적인\,재귀\,함수일\,때를\,중점으로\,두고\,만든\,클래스임을\,명심하고\,보는\,것이\,중요하다.
\\
\newline
\begin{lstlisting}[caption={재귀\,함수\,R의\,구동\,메소드\,Run}]
namespace RecursionFuntion
{
    public partial class RecursionFuntion<T,U>
    {
        public U Run(T x)
        {
            while (!g(x, f))
                x = f(x);
            return h(x, f, g);
        }
    }
}
\end{lstlisting}

위의\,내용을\,보면\,너무\,간단해서\,어이가\,없을\,수도\,있을\,것이다.\,
하지만\,위의\,내용은\,선형적인\,재귀(분기\,구조가\,없는\,재귀)에서\,나올\,수\,있는\,가장\,간단한\,식이며,\,기존의\,재귀\,함수에서도\,선형적인\,재귀\,함수는\,위와\,비슷하게\,$f(x)$의\,결과를\,x에\,햘당하고\,다시\,$f$에\,넣음으로서\,구동이\,가능하다.\,
이제\,재귀\,함수\,R을\,이용하여\,간단한\,양수\,팩토리얼\,재귀\,함수를\,만들어\,보자.
\\
\newline
\begin{lstlisting}[caption={재귀\,함수\,R을\,실제로\,사용하는\,예제}]
namespace RecursionFuntion
{
    internal class Program
    {
        public struct A
        {
            public int control { get; set; }
            public long value { get; set; }
            public A(int control, long value)
            {
                this.control = control;
                this.value = value;
            }

            public static A F(A x)
            {
                A a = new(x.control, x.value);
                a.value = a.value * a.control;
                a.control--;
                return a;
            }
            public static bool G(A x, Func<A,A> F)
            {
                if (x.control == 0 && F(x).value == 0)
                    return true;
                else
                    return false;
            }
            public static long H(A x, Func<A,A> F, Func<A,Func<A,A>,bool> G)
            {
                return x.value;
            }

        }
        static void Main(string[] args)
        {
            RecursionFuntion<A, long> recursionFuntion = new RecursionFuntion<A, long>(A.F, A.G, A.H);
            Console.WriteLine( recursionFuntion.Run(new(10,1)) );
        }
    }
}
\end{lstlisting}

위의\,내용을\,보면\,A.F에서\,연산을\,하고\,G가\,A와\,F를\,받아서\,조건을\,판단하고\,최종적으로\,H가\,해석해서\,내보내는\,구조임을\,알\,수\,있다.\,
이제\,다음으로\,Stack기반\,R을\,보러가자\\
다만\,Stack은\,초보여서\,구현이\,이상하게\,되어\,있을\,수\,있음으로\,코드\,전체와\,함수의\,정의에\,대해서\,말씀드리겠다.
\\
\newline
\begin{lstlisting}[caption={재귀\,함수\,R의\,Stack버전}]
namespace RecursionFuntion
{
    public class RecursionFuntionStack<T,U>
    {
        public Func<T, IEnumerable<T>> f { get; private set; }
        public Func<T, Func<T, IEnumerable<T>>, bool> g { get; private set; }
        public Func<T, Func<T, IEnumerable<T>>, Func<T, Func<T, IEnumerable<T>>, bool>, U> h { get; private set; }

        public RecursionFuntionStack() { }
        public RecursionFuntionStack(
            Func<T, IEnumerable<T>> f,
            Func<T, Func<T, IEnumerable<T>>, bool> g,
            Func<T, Func<T, IEnumerable<T>>, Func<T, Func<T, IEnumerable<T>>, bool>, U> h)
        {
            this.f = f;
            this.g = g;
            this.h = h;
        }

        public void SetF(Func<T, IEnumerable<T>> f)
        {
            this.f = f;
        }
        public void SetG(Func<T, Func<T, IEnumerable<T>>, bool> g)
        {
            this.g = g;
        }
        public void SetH(Func<T, Func<T, IEnumerable<T>>, Func<T, Func<T, IEnumerable<T>>, bool>, U> h)
        {
            this.h = h;
        }

        public List<U> Run(params T[] StartStates)
        {
            Stack<T> stack = new(StartStates);
            List<U> results = new();

            while (stack.Count > 0)
            {
                var current = stack.Pop();

                if (g(current,f))
                {
                    results.Add(h(current, f, g));
                }
                else
                {
                    foreach (var next in f(current))
                    {
                        stack.Push(next);
                    }
                }
            }
            return results;
        }

    }
}
\end{lstlisting}

여기에서\,$f,g,h$는\,다음과\,같이\,정의\,되어있다.\\
\newline
$\mathcal{T}$는\,$T$집합의\,요소인\,$T_n$들\,중에\,쓰는\,$T_n$의\,집합(배열이라고\,생각해도\,무방함)이다.\\
$\mathcal{U}$는\,$U$집합의\,요소인\,$U_n$들\,중에\,쓰는\,$U_n$의\,집합(배열이라고\,생각해도\,무방함)이다.\\
$f:T \rightarrow \mathcal{T}$\\
$g:T \times (T \rightarrow \mathcal{T}) \rightarrow \text{bool}$\\
$h:T \times (T \rightarrow \mathcal{T}) \times (T \times (T \rightarrow \mathcal{T}) \rightarrow \text{bool}) \rightarrow \mathcal{U}$\\
\newline
그러므로\,Stack버전\,재귀\,함수\,R은\,트리를\,다룰\,때에\,유용하게\,쓸\,수\,있을\,것으로\,추측하고\,실험삼아\,만든\,것이다.\\

그래서\,이\,구조는\,$R:T\rightarrow \mathcal{U}$인\,것이다.\,
그리고\,$h$가\,집합$\mathcal{U}$를\,내보내지\,않고\,최종적으로\,하나의\,$U$를\,내보내는\,경우도\,생각해볼\,수\,있다.\,
다만\,그렇게\,되도\,$h:T\times f \times g \rightarrow \mathcal{U}$임으로\,$\mathcal{U}$를\,$U$로\,축약하는\,해석기$h':\mathcal{U}\rightarrow U$를\,정의하면\,$R:T\rightarrow U$이면서\,비선형\,재귀인\,구조도\,가능할\,것이라고\,생각한다.\,
이때\,$R$은\,$R(f,g,h,h')(x)$로\,다시\,정의된다.
\begin{center}
$R(f,g,h,h')(x)=
\begin{cases}
h'(h(x,f,g)) &
\text{if}\,g(x,f(x))=\text{true} \\
R(f,g,h,h')(f(x)) &
\text{otherwise}
\end{cases}
$
\end{center}
\,
하지만\,$R:T\rightarrow U$인\,비선형\,재귀\,함수는\,아직\,구현하지\,않았다.
\\
\newline
\newline
이상으로\,문서를\,마치겠다.\,읽어주셔서\,고맙습니다!\\
\newline
-끝-

\end{document}